\documentclass{article}
\usepackage[margin=25mm]{geometry}
\usepackage{amsmath,amssymb}
\usepackage{xurl}
\usepackage[hidelinks]{hyperref}
\title{Finite Uniform Scope of the Matrix Bernstein Development}
\author{}
\date{October 28, 2026}
\begin{document}
\maketitle

\section{Source statement and present scope}

Tropp, \emph{User-friendly tail bounds for sums of random matrices},
Theorem~6.1(ii), gives a one-sided spectral Bernstein inequality for independent
Hermitian random matrices. Its rectangular version, Theorem~1.6, is the source
of the norm estimate quoted in Landau--Vazirani--Vidick, % chktex 8
arXiv:1307.5143v1, Theorem~2.13. See
\url{https://tropp.caltech.edu/papers/Tro11-User-Friendly-preprint.pdf} and
\url{https://arxiv.org/abs/1307.5143v1}.

The present theorem concerns independent uniform samples from one finite
nonempty family of Hermitian matrices. It is a specialization, rather than the
formalization of the source theorem for general matrix-valued random variables.
No general source-labelled theorem is marked proved by this result.

Let $Z:\mathcal A\to M_d(\mathbb C)$ be Hermitian, with
$\sum_{a\in\mathcal A}Z(a)=0$ and $\|Z(a)\|\le R$, where $R\ge0$.
Let $\sigma^2$ be any real upper bound for
\begin{align*}
	\left\|\frac1{|\mathcal A|}\sum_{a\in\mathcal A} Z(a)^{2}\right\|. % chktex 3
\end{align*}
For $m\ge0$ and $\tau\ge0$, the proved bound is
\begin{align*}
	\frac{\#\{\omega\in\mathcal A^m:
		\lambda_{\max}(\sum_j Z(\omega_j))\ge\tau\}}
	{|\mathcal A|^m}
	\le d\exp\left(-\frac{\tau^2}{2m\sigma^2+2R\tau/3}\right).
\end{align*}
The Lean event is expressed by the existence of a spectral value at least
$\tau$. The theorem includes zero variance and zero threshold. At a zero
denominator it uses the formal convention $0^{-1}=0$; the resulting bound by
$d$ is immediate. The coefficient $1/3$ and the variance norm are retained.

\section{Proof and extension plan}

Lieb's trace-log-exponential concavity is already proved on the full
positive-definite cone. Finite Jensen and induction over the sample coordinates
therefore give the matrix logarithmic moment generating function estimate.
The scalar bound
\begin{align*}
	e^x\le 1+x+\frac{x^2}{2(1-a/3)}
	\qquad (|x|\le a<3)
\end{align*}
follows from $2\,3^k\le{(k+2)!}$. Functional calculus, centering and the variance
norm bound give the exponential Markov estimate. Optimizing its scalar
parameter yields the displayed tail bound. These steps are derived in the
formal proof, rather than supplied as hypotheses.

The next finite step is the Hermitian dilation of an arbitrary complex matrix.
Its upper spectral tail gives the factor $2d$, and its squared diagonal blocks
give the maximum of the two Gram-matrix variance norms. Subsequently the finite
uniform expectations must be replaced by Bochner integrals and genuine
independence. The generic analytic and probability results belong in QICLean;
their transfer must precede any QICLean probability theorem depending on them.

The source-specific sampled projector is a further task. A norm approximation
alone does not establish the simultaneous bound on the frequency of each local
factor in every sampled word. The ground-energy dependence of the published
construction is recorded separately in
\path{docs/paper-gaps/lvv13_sampling_ground_energy.tex}.

\end{document}
